\documentclass[aspectratio=169]{beamer}
\usetheme{Copenhagen}
\usecolortheme{whale}
\usefonttheme[onlymath]{serif}

\usepackage{setspace}
\usepackage{xcolor}
\usepackage{graphicx}
\usepackage{multimedia}
\usepackage{fontawesome}
\usepackage[most]{tcolorbox}
\usepackage{empheq}
\usepackage{fancybox}
\usepackage{annotate-equations}
\usepackage{tikz}
\usetikzlibrary{positioning, calc, shapes.geometric, shapes.multipart, 
	shapes, arrows.meta, arrows, 
	decorations.markings, external, trees}
\tikzstyle{Arrow} = [
thick, 
decoration={
	markings,
	mark=at position 1 with {
		\arrow[thick]{latex}
	}
}, 
shorten >= 3pt, preaction = {decorate}
]

\definecolor{teel}{HTML}{00AEB3}
\definecolor{redorange}{HTML}{F26035}
\definecolor{darkred}{HTML}{8B0000}
\definecolor{darkblue}{HTML}{00008B}
\definecolor{darkgreen}{HTML}{06402B}
\definecolor{CarolinaBlue}{HTML}{4B9CD3}
\definecolor{deepsea}{HTML}{123456}
\definecolor{sharedblue}{HTML}{009292}
\definecolor{distalblue}{HTML}{490092}

\usecolortheme[named=CarolinaBlue]{structure}


\title[Academia]{Faculty Positions in the United States}
\author[Paul Zivich]{Zivich Research Group}
\date{~}

\setbeamercovered{transparent}
\setbeamertemplate{navigation symbols}{}  % gets rid of the dumb navigation symbols
\setbeamertemplate{page number in head/foot}{\insertframenumber}  % adds slide #
\setbeamertemplate{headline}{}

\AtBeginSection[]{
	\begin{frame}
		\vfill
		\centering
		\begin{beamercolorbox}[sep=8pt,center,shadow=true,rounded=true]{title}
			\usebeamerfont{title}\insertsectionhead\par%
		\end{beamercolorbox}
		\vfill
	\end{frame}
}

\begin{document}

\maketitle

\begin{frame}{Faculty Positions}
	Describe basic structure of academic positions in the United States
	~\\~\\
	Help you to understand structures and career path differences
\end{frame}

\begin{frame}{Academic Tracks}
	Different institutions have different tracks (or names). At UNC,
	\begin{itemize}
		\item \textbf{Tenure}: what you usually think of
		\item \textbf{Research}: primary role is research
		\item \textbf{Clinical}: have clinical duties outside of other roles
	\end{itemize}
	~\\
	Each role entails different responsibilities
\end{frame}

\begin{frame}{Academic Tracks}
	Differences:
	\begin{itemize}
		\item \textbf{Tenure}: research, teaching (expected to teach a set amount of credit hours), service (part of academic committees in department, school, university)
		\item \textbf{Research}: primarily research with fewer obligations to department and university, but expected to have full external coverage of your effort
	\end{itemize}
	~\\
	The following positions are all on one of these tracks
	\begin{itemize}
		\item This determines roles and promotion criteria
	\end{itemize}
\end{frame}

\begin{frame}{Assistant Professor}
	Entry-level professor position
	\begin{itemize}
		\item First step on the ladder
		\item Occurs directly after graduation or after a postdoctoral period
	\end{itemize}
\end{frame}

\begin{frame}{Associate Professor}
	Mid-level professor position
	\begin{itemize}
		\item Next promotion level in academia
		\item Sometimes jointly given with tenure
		\item More responsibilities in department
	\end{itemize}
	~\\
	Tenure
	\begin{itemize}
		\item A more permanent position in academia. Quite difficult for a university to dismiss someone who has received tenure
	\end{itemize}
\end{frame}

\begin{frame}{Professor}
	High-level professor position
	\begin{itemize}
		\item Highest-level generally
		\item Some variations at this level
	\end{itemize}
	~\\
	Full Professor
	\begin{itemize}
		\item Have voting privileges on department policy and decisions
	\end{itemize}
\end{frame}

\begin{frame}{Others}
	Adjunct Professor
	\begin{itemize}
		\item Part-time, usually teaching
	\end{itemize}
	~\\
	Distinguished Professor
	\begin{itemize}
		\item Usually have set salary coverage, named by a donor
	\end{itemize}
	~\\
	Professor Emeritus
	\begin{itemize}
		\item Retired, with some access to university resources
	\end{itemize}	
\end{frame}

\begin{frame}{Questions?}
	~
\end{frame}

\end{document}
